% Figure (e) v4: the checker's unresolved window on a real gem5 trace.
% Body only (goes into the paper's tikz/). Everything is wrapped in { } so no
% macro or style leaks into other figures; all helper names are prefixed \cw.
% Needs TikZ libraries calc + shapes.geometric (the paper preamble loads both).
% Data: STAGE3_gem5/tests/fixtures/inv_uv1.pipeview.txt.gz (STT_6
% mispredict_not_taken/inst-002014, amulet InvisiSpec, --ruby
% --scheme=SpectreSafeInvisibleSpec, resolve stall 5000). Cycle c = (tick - 3157000)/500.
% Every number, label and source: paper_figures/plans/fig_e_checkerwindow.md.
% Text convention (see the plan): sentence case for prose; lc, fnc, xmit are fixed
% lowercase abbreviations set with \cwab (one font everywhere, upright even in italics).
{%
% ============================ TUNING KNOBS ============================
\def\cwcyc{6.5}       % pt per simulated cycle (same scale in every axis segment)
\def\cwbrk{14}        % pt width of each elided stretch (dotted)
\def\cwtl{66}         % pt, x of the time axis origin (cycle -0.5); left of it: label column
\def\cwrowh{30}       % pt between instruction rows
\def\cwbarh{4}        % pt height of an instruction bar
\def\cwbusx{1.5}      % pt, x of the Alloy-label bus line
\def\cwlabx{6}        % pt, left edge of the row labels
\def\cwiconx{3.85}    % pt, centre x of the source icons
\def\cwicongap{11}    % pt, left edge of a source icon's label
\def\cwright{250}     % pt, right edge of the criterion box (figure is <= 252pt wide)
\def\cwwinlw{0.9pt}   % line width of the window frame
\def\cwstar{8.5pt}    % outer diameter of the transmission star
\def\cwwtop{13}       % pt, window frame top above row i0
\def\cwwbot{25}       % pt, window frame bottom below row i2
\def\cwgap{10}        % pt between the cycle numbers and the criterion box
\def\cwdefshade{black!12} % fill marking defense-specific elements (the only gray)
\def\cwfs{\scriptsize}    % the one font size (7pt)
% axis segments [start,end], in cycles after i0's fetch; elided stretches between them
\def\cwsA{-0.5}\def\cweA{8.5}       % front end; lc retires at 8
\def\cwsB{39.5}\def\cweB{41.5}      % xmit issues at 40
\def\cwsC{68.5}\def\cweC{71.5}      % transmission (L1D eviction) at 70
\def\cwsE{5006.5}\def\cweE{5013.5}  % branch resolves 5007, squash 5008, fnc retires 5013
% ======================================================================
\pgfmathsetmacro\cwoA{0}%
\pgfmathsetmacro\cwoB{(\cweA-\cwsA)*\cwcyc+\cwbrk}%
\pgfmathsetmacro\cwoC{\cwoB+(\cweB-\cwsB)*\cwcyc+\cwbrk}%
\pgfmathsetmacro\cwoE{\cwoC+(\cweC-\cwsC)*\cwcyc+\cwbrk}%
\pgfmathsetmacro\cwend{\cwtl+\cwoE+(\cweE-\cwsE)*\cwcyc}%
% x (pt) of cycle #2 in segment #1 (A, B, C or E)
\newcommand\cwx[2]{(\cwtl+\csname cwo#1\endcsname+(#2-\csname cws#1\endcsname)*\cwcyc)}%
% fixed abbreviation (lc, fnc, xmit): always upright roman, never capitalized
\newcommand\cwab[1]{\textup{#1}}%
% bar edges from x=#2 to x=#3 at row y=#1; #4 = extra draw options
\newcommand\cwedges[4]{\draw[#4] (#2,#1+\cwbarh/2) -- (#3,#1+\cwbarh/2) (#2,#1-\cwbarh/2) -- (#3,#1-\cwbarh/2);}%
% vertical tick across a bar at x=#2, row y=#1
\newcommand\cwtck[2]{\draw (#2,#1-\cwbarh/2) -- (#2,#1+\cwbarh/2);}%
% stage letter #3 above the bar at x=#2, row y=#1
\newcommand\cwstg[3]{\node[anchor=south] at (#2,#1+\cwbarh/2+1.6) {#3};}%
% bar from x=#2 (in segment A) through every elided stretch to x=#3 (in segment E), row y=#1
\newcommand\cwlongbar[3]{%
  \cwedges{#1}{#2}{\cwtl+\cwoB-\cwbrk}{}%
  \cwedges{#1}{\cwtl+\cwoB-\cwbrk}{\cwtl+\cwoB}{densely dotted}%
  \cwedges{#1}{\cwtl+\cwoB}{\cwtl+\cwoC-\cwbrk}{}%
  \cwedges{#1}{\cwtl+\cwoC-\cwbrk}{\cwtl+\cwoC}{densely dotted}%
  \cwedges{#1}{\cwtl+\cwoC}{\cwtl+\cwoE-\cwbrk}{}%
  \cwedges{#1}{\cwtl+\cwoE-\cwbrk}{\cwtl+\cwoE}{densely dotted}%
  \cwedges{#1}{\cwtl+\cwoE}{#3}{}}%
% document icon centred at (#1,#2), fill #3, glyph #4 (empty = text lines); 7.7 x 10 pt
\newcommand\cwdoc[4]{%
  \begin{scope}[shift={(#1,#2)}]
    \draw[fill=#3] (-3.85,-5) -- (3.85,-5) -- (3.85,2.95) -- (1.8,5) -- (-3.85,5) -- cycle;
    \draw (1.8,5) -- (1.8,2.95) -- (3.85,2.95);
    \if\relax\detokenize{#4}\relax
      \foreach \cwyy in {1.4,-0.7,-2.8} \draw[line width=0.3pt] (-2.3,\cwyy) -- (2.3,\cwyy);
    \else \node[font=\tiny] at (0,-0.9) {#4};\fi
  \end{scope}}%
\begin{tikzpicture}[x=1pt, y=1pt, >=latex,
  every node/.style={font=\cwfs, inner sep=0pt,
    execute at begin node={\hyphenpenalty=10000\exhyphenpenalty=10000\relax}},
  cwann/.style={font=\cwfs\itshape},
  cwrow/.style={anchor=base west},
  cwgloss/.style={font=\cwfs\itshape, anchor=base west},
]
\pgfmathsetmacro\cwyA{0}
\pgfmathsetmacro\cwyB{-\cwrowh}
\pgfmathsetmacro\cwyC{-2*\cwrowh}
\pgfmathsetmacro\cwwT{\cwyA+\cwwtop}      % window frame top
\pgfmathsetmacro\cwwB{\cwyC-\cwwbot}      % window frame bottom
\pgfmathsetmacro\cwyT{\cwwB-3.5}          % top of the cycle numbers
\pgfmathsetmacro\cwyK{\cwyT-7-\cwgap}     % top of the criterion box
\pgfmathsetmacro\cwyS{\cwwT+7.5}          % centre line of the window title / Alloy labels
\pgfmathsetmacro\cwyL{\cwyS+13}           % legend line

% ------------------------------ legend line ------------------------------
\node[anchor=west] at (0,\cwyL) {F: fetch\quad D: decode\quad E: issue\quad Dotted: elided cycles};
\node[anchor=east] (cwlg) at (\cwright,\cwyL) {Defense-specific};
\draw[fill=\cwdefshade] ($(cwlg.west)+(-9.5,-3)$) rectangle ++(6,6);

% ------------- left column: the roles come from the Alloy instance -------------
\cwdoc{\cwiconx}{\cwyS}{white}{\{\,\}}
\node[anchor=base west] at (\cwicongap,\cwyS-2.4) {Role labels from the Alloy instance};
\draw (\cwbusx,\cwyS-5) -- (\cwbusx,\cwyC);
\foreach \cwy in {\cwyA,\cwyB,\cwyC} \draw (\cwbusx,\cwy) -- (\cwlabx-2,\cwy);
% the id sits level with its bar; the gloss goes underneath
% (one 8pt baseline grid: id baseline = bar centre - 2.4pt, glosses 8pt and 16pt lower)
\node[cwrow] at (\cwlabx,\cwyA-2.4) {i0: \cwab{lc}};
\node[cwgloss] at (\cwlabx,\cwyA-10.4) {Last committed};
\node[cwrow] at (\cwlabx,\cwyB-2.4) {i1: \cwab{fnc}};
\node[cwgloss] at (\cwlabx,\cwyB-10.4) {First non-committed};
\node[cwgloss] at (\cwlabx,\cwyB-18.4) {(unresolved branch)};
\node[cwrow] at (\cwlabx,\cwyC-2.4) {i2: \cwab{xmit}};
\node[cwgloss] at (\cwlabx,\cwyC-10.4) {Transmitter};

% ------------------- i0: last committed (retires at 8) -------------------
\cwedges{\cwyA}{{\cwx{A}{0}}}{{\cwx{A}{8}}}{}
\cwtck{\cwyA}{{\cwx{A}{0}}} \cwtck{\cwyA}{{\cwx{A}{1}}} \cwtck{\cwyA}{{\cwx{A}{4}}} \cwtck{\cwyA}{{\cwx{A}{8}}}
\cwstg{\cwyA}{{\cwx{A}{0}}}{F} \cwstg{\cwyA}{{\cwx{A}{1}}}{D} \cwstg{\cwyA}{{\cwx{A}{4}}}{E}
\node[cwann, anchor=west] at ({\cwx{A}{8}+3.5},\cwyA) {Retires (window opens)};

% ------------- i1: first non-committed = the unresolved branch -------------
\cwlongbar{\cwyB}{{\cwx{A}{1}}}{{\cwx{E}{5013}}}
\cwtck{\cwyB}{{\cwx{A}{1}}} \cwtck{\cwyB}{{\cwx{A}{2}}} \cwtck{\cwyB}{{\cwx{A}{6}}} \cwtck{\cwyB}{{\cwx{E}{5007}}} \cwtck{\cwyB}{{\cwx{E}{5013}}}
\cwstg{\cwyB}{{\cwx{A}{1}}}{F} \cwstg{\cwyB}{{\cwx{A}{2}}}{D} \cwstg{\cwyB}{{\cwx{A}{6}}}{E}
\node[cwann, anchor=south] at ({(\cwx{A}{8}+\cwx{E}{5007})/2},\cwyB+\cwbarh/2+2) {5000-cycle resolve stall};
\node[cwann, anchor=north east] at ({\cwx{E}{5007}-3},\cwyB-\cwbarh/2-2.5) {Resolves (window closes)};
\node[cwann, anchor=south east] at ({\cwx{E}{5013}+0.5},\cwyB+\cwbarh/2+2) {Retires};

% ------------------- i2: the transmitter (a load), squashed -------------------
\cwlongbar{\cwyC}{{\cwx{A}{2}}}{{\cwx{E}{5008}}}
\cwtck{\cwyC}{{\cwx{A}{2}}} \cwtck{\cwyC}{{\cwx{A}{3}}} \cwtck{\cwyC}{{\cwx{B}{40}}}
\cwstg{\cwyC}{{\cwx{A}{2}}}{F} \cwstg{\cwyC}{{\cwx{A}{3}}}{D} \cwstg{\cwyC}{{\cwx{B}{40}}}{E}
\draw[line width=0.6pt] ({\cwx{E}{5008}-2},\cwyC-2) -- ++(4,4) ({\cwx{E}{5008}-2},\cwyC+2) -- ++(4,-4);
\node[cwann, anchor=west] at ({\cwx{E}{5008}+4.5},\cwyC) {Squashed};
% the transmission: the one star in the figure
\node[star, star points=5, star point ratio=2.4, fill=black, draw=white, line width=0.5pt,
      minimum size=\cwstar, inner sep=0pt] (cwtx) at ({\cwx{C}{70}},\cwyC) {};
\node[anchor=north, align=center] at ($(cwtx.south)+(0,-2.5)$) {Transmission\\[1pt]\textit{(L1D eviction)}};

% ------------------------- the unresolved window -------------------------
\draw[line width=\cwwinlw] ({\cwx{A}{8}},\cwwT) rectangle ({\cwx{E}{5007}},\cwwB);
\node[anchor=base] at ({(\cwx{A}{8}+\cwx{E}{5007})/2},\cwyS-2.4) {Unresolved window};

% ------------------------------ time axis ------------------------------
\node[anchor=north west] at (0,\cwyT) {Cycle};
\foreach \cwc/\cws in {0/A, 8/A, 40/B, 70/C, 5007/E, 5013/E}
  \node[anchor=north] at ({\cwx{\cws}{\cwc}},\cwyT) {\cwc};

% ------ defense-specific transmission criterion, checked inside the window ------
\pgfmathsetmacro\cwcw{\cwright-\cwtl-6}
\node[draw, fill=\cwdefshade, anchor=north west, align=left, inner sep=3pt,
      text width=\cwcw pt] (cwcrit) at (\cwtl,\cwyK)
  {\raggedright\textit{InvisiSpec, checked inside the window:} a fill, eviction, invalidation, coherence or replacement update in the L1D, L2, directory or D-TLB by a squashed instruction younger than \cwab{lc}. Speculative-buffer fills do not count.};
\draw[->] ({(\cwx{B}{40}+\cwx{C}{70})/2},\cwyK) -- ({(\cwx{B}{40}+\cwx{C}{70})/2},\cwwB);
\path (cwcrit.center); \pgfgetlastxy{\cwqx}{\cwqy}
\cwdoc{\cwiconx}{\cwqy}{\cwdefshade}{}
\node[anchor=west, align=left] (cwcl) at (\cwicongap,\cwqy) {Defense\\transmission\\criterion};
\draw[->] ($(cwcl.east |- cwcrit.west)+(2,0)$) -- (cwcrit.west);

\end{tikzpicture}%
}%
